\documentclass{article}
\usepackage{amsmath}
\begin{document}
\section*{Controller Target Tracking Audit}
At rollout time $t_i$, the controller target $r_j(t_i)$ is the linear
interpolant of the recorded position track for coordinate $j$ on the same
relative time axis. The reported sampled residual is $e_{ij}=q_{ij}-r_j(t_i)$.
For $N$ samples, the reported values are
\[
 R_j=\sqrt{N^{-1}\sum_{i=1}^{N} e_{ij}^2},\qquad
 M_j=\max_i |e_{ij}|.
\]
The velocity reference is the source controller's finite-difference rate
table, differentiated separately on each side of the configured impact and
sampled with the same side-specific interpolation and endpoint clamps.
Velocity errors use the same sampled RMS and maximum formulas in native
coordinate units per second. These derived rates are not measured velocities
and are not derivatives of the linearly interpolated position target.
Coordinates keep their native units and are reported in model order; no mixed
translation--rotation aggregate is defined. Only finite, increasing and
covered time axes are accepted, except that the final state may exceed the
reference endpoint by at most half the fixed 1 ms step, when rounding requires
it. The controller clamps that target to its last value; the receipt reports
the clamped duration. The record digest must agree with the bundle
provenance; the run model with its recorded shaft configuration must match the
bundle model. This comparison does not establish that the record was generated
from that model, data provenance, correct model
migration, continuity between samples, actuator feasibility or impact.
For a settled lead-in of $L$ fixed steps, the run repeats those steps, the
score begins at step $L$, and the controller target is sampled at the elapsed
time $t_i+L\Delta t$. Only the existing no-impact fixed-step lead-in contract
is supported; inconsistent metadata is rejected. This does not establish
independent provenance for the settled state.
\end{document}
